% Самостоятельная методичка лекции 9. Сборка из корня:
% python3 9/images/geometry_probability.py
% python3 9/images/avoidance_example.py
% latexmk 9/collision-avoidance.tex
\documentclass[a4paper,11pt]{article}
\usepackage{iftex}
\RequirePDFTeX
\usepackage{cmap}
\usepackage[T2A]{fontenc}
\usepackage[utf8]{inputenc}
\usepackage[english,russian]{babel}
\usepackage{PTSerif}
\usepackage{amsmath,amssymb,amsthm,mathtools}
\usepackage[a4paper,left=20mm,right=20mm,top=22mm,bottom=23mm,headheight=14pt]{geometry}
\usepackage{microtype}
\usepackage{graphicx,xcolor,tikz,booktabs,tabularx,array,longtable}
\usetikzlibrary{arrows.meta,calc,positioning}
\usepackage{enumitem,fancyhdr,needspace,float,placeins}
\usepackage[most]{tcolorbox}
\definecolor{accent}{HTML}{187AA5}
\definecolor{muted}{HTML}{5A6166}
\definecolor{figink}{HTML}{25323F}
\definecolor{figtan}{HTML}{087F78}
\definecolor{figorange}{HTML}{B57714}
\definecolor{figpurple}{HTML}{79539C}
\newtcolorbox{explanation}[1]{enhanced,
 colback=accent!3!white,colframe=accent!80!black,
 boxrule=.5pt,toprule=1.2pt,arc=1mm,
 left=3mm,right=3mm,top=3mm,bottom=2mm,
 before skip=9pt,after skip=9pt,
 title={#1},fonttitle=\bfseries,coltitle=accent!70!black,
 attach boxed title to top left={xshift=3mm,yshift=-2mm},
 boxed title style={colback=white,frame hidden,boxrule=0pt,
 left=1.5mm,right=1.5mm,top=.4mm,bottom=.4mm}}
\newtcolorbox{remarkbox}{enhanced,
 colback=orange!4!white,colframe=orange!65!black,
 boxrule=.5pt,toprule=1.2pt,arc=1mm,
 left=3mm,right=3mm,top=2mm,bottom=2mm,
 before skip=9pt,after skip=9pt}
\usepackage{hyperref}
\usepackage{bookmark}
\hypersetup{unicode=true,colorlinks=true,linkcolor=accent,citecolor=accent,urlcolor=accent,
 pdftitle={Манёвры уклонения от столкновения},
 pdfauthor={Материалы курса «Введение в маневрирование КА»},
 pdfsubject={Геометрия сближения, вероятность столкновения и выбор импульсного манёвра}}
\newcommand{\dd}{\mathop{}\!\mathrm d}
\newcommand{\vect}[1]{\boldsymbol{#1}}
\newcommand{\norm}[1]{\left\lVert#1\right\rVert}
\DeclareMathOperator{\sinc}{sinc}
\numberwithin{equation}{section}
\setcounter{tocdepth}{2}
\setcounter{secnumdepth}{2}
\setlength{\parindent}{1.3em}
\setlength{\parskip}{2pt}
\linespread{1.055}
\setlist{itemsep=2pt,topsep=4pt,parsep=0pt}
\setlength{\emergencystretch}{1.5em}
\clubpenalty=10000
\widowpenalty=10000
\displaywidowpenalty=10000
\allowdisplaybreaks[1]
\setcounter{topnumber}{2}
\renewcommand{\topfraction}{.88}
\renewcommand{\textfraction}{.10}
\renewcommand{\floatpagefraction}{.82}
\pagestyle{fancy}
\fancyhf{}
\fancyhead[L]{\small\color{muted}Манёвры уклонения от столкновения}
\fancyfoot[C]{\small\thepage}
\renewcommand{\headrulewidth}{.3pt}
\newcommand{\coursefigure}[3][.98]{%
\begin{figure}[htbp]\centering
\includegraphics[width=#1\textwidth]{images/#2.pdf}
\caption{#3}\label{fig:#2}\end{figure}}
\DeclareMathOperator{\diag}{diag}
\DeclareMathOperator{\rank}{rank}
% Generated by avoidance_example.py
\newcommand{\ExamplePeriod}{5828{,}517}
\newcommand{\ExampleLead}{2331{,}407}
\newcommand{\ExampleNominalMiss}{29{,}155}
\newcommand{\ExampleNominalPc}{5{,}88\cdot10^{-3}}
\newcommand{\ExampleRelativeSpeed}{7{,}5842}
\newcommand{\ExampleOptimalDV}{0{,}097840}
\newcommand{\ExampleOptimalMiss}{348{,}632}
\newcommand{\ExampleOptimalNonlinearPc}{1{,}00\cdot10^{-5}}
\newcommand{\ExampleOptimalNonlinearMiss}{348{,}637}
\newcommand{\ExampleOptimalTCA}{0{,}026143}
\newcommand{\ExampleExecutionPc}{1{,}18\cdot10^{-5}}
\newcommand{\ExampleSafeDV}{0{,}100996}
\newcommand{\ExampleSafePc}{8{,}00\cdot10^{-6}}
\newcommand{\ExampleSafeNonlinearPc}{8{,}00\cdot10^{-6}}
\newcommand{\ExampleSafeNonlinearMiss}{360{,}573}
\newcommand{\ExampleSafeTCA}{0{,}026988}
\newcommand{\ExampleGRepeat}{1{,}22\cdot10^{-9}}
\newcommand{\ExamplePcRepeat}{6{,}60\cdot10^{-15}}
\newcommand{\ExampleMatrixG}{\begin{pmatrix}-483{,}069&-3321{,}815&0{,}000\\2882{,}996&4276{,}692&-256{,}984\end{pmatrix}}
\newcommand{\ExampleCovarianceMatrix}{\begin{pmatrix}14233{,}213&-18042{,}098\\-18042{,}098&27366{,}787\end{pmatrix}}

\begin{document}
\begin{titlepage}
\centering
{\large Московский физико-технический институт\par}
\vspace{8mm}
{\large Курс «Введение в маневрирование КА»\par}
\vspace{32mm}
{\color{accent}\rule{28mm}{.8pt}\par}
\vspace{9mm}
{\LARGE\bfseries Манёвры уклонения\\[3mm]
от столкновения\par}
\vspace{11mm}
{\large Методические материалы к лекции 9\par}
\vspace{8mm}
{Геометрия сближения и неопределённость.\\
Вероятность столкновения и выбор импульса.\\
Численный расчёт и проверка манёвра.\par}
\vfill
{Рабочая редакция\par}
\vspace{3mm}
{13 ноября 2026 года\par}
\end{titlepage}
\setcounter{page}{2}
\begingroup\small
\setlength{\parskip}{0pt}
\tableofcontents
\endgroup
\clearpage

\section*{Предмет рассмотрения и обозначения}
\addcontentsline{toc}{section}{Предмет рассмотрения и обозначения}

Уклонение изменяет движение космического аппарата так, чтобы снизить
опасность прогнозируемого сближения. Для выбора манёвра недостаточно
знать только номинальное расстояние между объектами: необходимо
учитывать неопределённость их положения и размеры. Поэтому расчёт
связывает три задачи: прогноз движения, оценку вероятности столкновения
и выбор допустимого управления.

Рассматривается короткое быстрое сближение управляемого аппарата
с объектом, движение которого не изменяется манёвром. Сначала
выводится двумерная модель вероятности столкновения, затем строится
линейная связь импульса со смещением в плоскости сближения.
Эта связь продолжает использованную в задаче встречи запись
$M=\Phi B$. Основной расчёт проводится для пространственных
эллиптических орбит в центральном поле; выбранный импульс проверяется
повторным нелинейным прогнозом.

\begin{center}
\renewcommand{\arraystretch}{1.12}
\begin{tabularx}{\textwidth}{@{}>{$}l<{$}X@{}}
\toprule
\multicolumn{1}{@{}l}{Обозначение}&Смысл\\\midrule
\vect R_k,\vect V_k & Инерциальные положение и скорость объекта $k$;
$k=1$ --- управляемый КА, $k=2$ --- опасный объект.\\
\vect y & Классические элементы $(a,e,\omega,i,\Omega,M)^\top$.\\
\vect\rho,\vect v_{\mathrm{rel}} & Относительные положение
$\vect R_2-\vect R_1$ и скорость $\vect V_2-\vect V_1$ в инерциальных осях.\\
t_{\mathrm{ca}},t_m & Момент наибольшего сближения и момент импульсного манёвра.\\
P,\vect b & Матрица проекции на плоскость сближения и двумерный вектор промаха.\\
R_c & Сумма радиусов сфер, описанных около двух объектов.\\
C_{\rho},C_b & Ковариации относительного положения в пространстве и в плоскости сближения.\\
P_c,P_{\text{доп}} & Вероятность столкновения за рассматриваемый пролёт и её допустимый уровень.\\
\Delta\vect V & Импульс скорости в предманёвренной ОСК;
компоненты $(\Delta V_r,\Delta V_t,\Delta V_n)^\top$.\\
B,\Phi,M & Матрица чувствительности элементов, матрица изохронных производных
и матрица манёвра $M=\Phi B$.\\
G & Матрица влияния импульса на вектор промаха: $\delta\vect b=G\Delta\vect V$.\\
J & Характеристическая скорость; для одного импульса $J=\norm{\Delta\vect V}$.\\
\bottomrule
\end{tabularx}
\end{center}
Углы в формулах выражены в радианах. В формулах одной задачи
используются согласованные единицы длины, скорости и ковариаций.
Эксцентрическая аномалия обозначается $E$,
математическое ожидание --- $\mathbb{E}$.

% Главы 1--3 включаются сюда при окончательной сборке.
\section{Постановка задачи уклонения}
\label{sec:formulation}

При переходе требуется получить заданную орбиту, при встрече ---
заданное состояние в установленный момент. В задаче уклонения
конечное состояние обычно не фиксируется. Манёвр должен изменить
ожидаемое сближение так, чтобы вероятность столкновения стала
допустимой, а нарушение основной программы полёта оставалось
приемлемым. Поэтому расчёт начинается с определения опасного события
и меры его опасности.

\subsection{Исходные данные и управление}

Рассмотрим управляемый космический аппарат, обозначенный индексом~1,
и второй объект, обозначенный индексом~2. На начальную эпоху заданы
их состояния и ковариационные матрицы ошибок определения состояний.
Выбранная модель движения позволяет прогнозировать оба объекта
на интервале, содержащем ожидаемое сближение. В нашей расчётной
модели второй объект не маневрирует.

Движение каждого объекта можно описывать классическими элементами
\[
\vect y=(a,e,\omega,i,\Omega,M)^\top.
\]
Для оценки сближения эти элементы преобразуются в инерциальные
положение $\vect R$ и скорость $\vect V$. Геометрические сравнения
выполняются при одном физическом времени и в одной системе координат.
Это относится и к номинальным состояниям, и к их ковариациям.

Сначала ограничимся одним импульсным манёвром. В момент $t_m$
положение аппарата непрерывно, а его инерциальная скорость получает
приращение:
\begin{equation}
\vect R_1(t_m^+)=\vect R_1(t_m^-),\qquad
\vect V_1(t_m^+)=\vect V_1(t_m^-)+\Delta\vect V_{\text{ИСК}}.
\label{eq:initial-impulse-jump}
\end{equation}
Для задания управления удобно использовать компоненты импульса
в орбитальной системе координат аппарата в момент манёвра:
\begin{equation}
\Delta\vect V=
\begin{pmatrix}\Delta V_r&\Delta V_t&\Delta V_n\end{pmatrix}^{\!\top},
\qquad
\Delta\vect V_{\text{ИСК}}=
\begin{bmatrix}\vect e_r&\vect e_t&\vect e_n\end{bmatrix}
\Delta\vect V.
\label{eq:rtn-impulse}
\end{equation}
Базис ортонормирован, поэтому модуль импульса не зависит от выбора
этих двух систем координат. После скачка скорости вычисляется новая
орбита и продолжается прогноз движения.

Искомыми параметрами являются время $t_m$, направление и величина
импульса. Время ограничено снизу сроком подготовки манёвра,
сверху --- временем, после которого требуемое уклонение уже
невыполнимо. Допустимые направления определяются ориентацией
аппарата и расположением двигателей. Ограничение величины импульса
связано с ресурсом двигательной установки и допустимым изменением
рабочей орбиты.

\subsection{Критерий и условие безопасности}

Пусть $P_c(t_m,\Delta\vect V)$ --- оценённая вероятность столкновения
в рассматриваемом сближении после манёвра, а $P_{\text{доп}}$ ---
принятое для данной задачи допустимое значение. Тогда простая
задача минимизации характеристической скорости имеет вид
\begin{equation}
\begin{aligned}
\min_{\mathclap{t_m,\,\Delta\vect V\in\mathcal A(t_m)}}\quad
 &\norm{\Delta\vect V},\\
\text{s.t.}\quad
 &P_c(t_m,\Delta\vect V)\leqslant P_{\text{доп}},\\
 &t_m\in[t_{m,\min},t_{m,\max}].
\end{aligned}
\label{eq:cam-problem}
\end{equation}
Множество $\mathcal A(t_m)$ задаёт допустимые направления
и величины импульса в выбранный момент манёвра.
Для нескольких импульсов критерий заменяется суммой их модулей.
В последующих разделах подробно рассматривается одноимпульсная
задача: она позволяет проследить связь направления манёвра,
геометрии сближения и вероятности столкновения.

Условие на $P_c$ относится к конкретной вероятностной модели:
оно зависит от принятой неопределённости состояний, размеров
объектов и способа расчёта. Допустимый порог является условием
задачи. Он не следует из уравнений движения.

\begin{explanation}{Что требуется определить}
Расчёт уклонения должен дать манёвр, удовлетворяющий условию
безопасности и ограничениям миссии. Увеличение номинального
расстояния является промежуточным результатом. Для проверки
условия~\eqref{eq:cam-problem} необходимо учесть также неопределённость
взаимного положения объектов.
\end{explanation}

\section{Геометрия сближения}
\label{sec:geometry}

До выбора управления необходимо определить, когда объекты
окажутся ближе всего и в каких направлениях изменение положения
влияет на промах. Для этого сначала рассмотрим номинальные
траектории. Неопределённость будет добавлена в следующей главе.

\subsection{Момент наибольшего сближения}

В инерциальной системе определим относительное положение и скорость:
\begin{equation}
\vect\rho(t)=\vect R_2(t)-\vect R_1(t),\qquad
\vect v_{\mathrm{rel}}(t)=\vect V_2(t)-\vect V_1(t).
\label{eq:relative-inertial}
\end{equation}
Здесь $\dot{\vect\rho}=\vect v_{\mathrm{rel}}$. Производная относительных
координат во вращающейся орбитальной системе отличается от этой
скорости на член, связанный с вращением базиса.

Момент наибольшего сближения $t_{\mathrm{ca}}$ (time of closest
approach, TCA) минимизирует $\norm{\vect\rho(t)}$ в рассматриваемой
окрестности события. Для внутреннего минимума
\begin{equation}
\frac{\dd}{\dd t}\frac{\norm{\vect\rho}^2}{2}
=\vect\rho^\top\vect v_{\mathrm{rel}}=0.
\label{eq:tca-condition}
\end{equation}
То есть в TCA вектор относительного положения перпендикулярен
относительной скорости. Для невырожденного локального минимума дополнительно
требуется
\begin{equation}
\norm{\vect v_{\mathrm{rel}}}^2+
\vect\rho^\top(\vect a_2-\vect a_1)>0.
\label{eq:tca-curvature}
\end{equation}
Номинальное расстояние в этом событии равно
$d_{\min}=\norm{\vect\rho(t_{\mathrm{ca}})}$.
Если на интервале прогноза найдено несколько локальных минимумов,
каждое соответствующее сближение оценивается отдельно.

Для короткого пролёта вблизи эпохи $t_0$ можно пренебречь изменением
относительной скорости:
\begin{equation}
\vect\rho(t_0+\tau)\approx
\vect\rho(t_0)+\vect v_{\mathrm{rel}}(t_0)\tau.
\label{eq:straight-line}
\end{equation}
После подстановки в квадрат расстояния получаем
\[
\norm{\vect\rho(t_0)+\vect v_{\mathrm{rel}}\tau}^2
=\norm{\vect\rho(t_0)}^2+
2\tau\vect\rho(t_0)^\top\vect v_{\mathrm{rel}}
+\tau^2\norm{\vect v_{\mathrm{rel}}}^2.
\]
Дифференцирование по $\tau$ даёт оценку поправки ко времени:
\begin{equation}
t_{\mathrm{ca}}-t_0\approx
-\frac{\vect\rho(t_0)^\top\vect v_{\mathrm{rel}}(t_0)}
       {\norm{\vect v_{\mathrm{rel}}(t_0)}^2}.
\label{eq:linear-tca}
\end{equation}
Эта формула пригодна для локального уточнения уже найденного
быстрого сближения. После манёвра меняются обе части
условия~\eqref{eq:tca-condition}; поэтому прежняя эпоха TCA
обычно перестаёт быть моментом минимального расстояния.

\subsection{Плоскость сближения}

Пусть относительная скорость в номинальной TCA ненулевая.
Введём единичный вектор
$\vect e_v=\vect v_{\mathrm{rel}}/\norm{\vect v_{\mathrm{rel}}}$
и два ортонормированных вектора $\vect e_1,\vect e_2$,
перпендикулярных $\vect e_v$. Они задают \emph{плоскость сближения}.
В литературе по манёврам уклонения используется также название
$b$-плоскость~\cite{gonzalo}.
Проекция на неё записывается матрицей
\begin{equation}
P=\begin{bmatrix}\vect e_1^\top\\\vect e_2^\top\end{bmatrix},
\qquad
PP^\top=I_2,\qquad
P^\top P=I_3-\vect e_v\vect e_v^\top.
\label{eq:encounter-projection}
\end{equation}
Вектор промаха в координатах плоскости равен
\begin{equation}
\vect b=P\vect\rho.
\label{eq:miss-plane}
\end{equation}
В прямолинейной модели $P\vect v_{\mathrm{rel}}=0$, поэтому
$P\vect\rho(t)$ не изменяется при движении вдоль прямой.
Продольное перемещение изменяет время прохождения сближения,
а поперечное --- промах. Именно это разделение делает плоскость
удобной для расчёта уклонения.

В точной TCA выполнено $\vect\rho\perp\vect v_{\mathrm{rel}}$,
и потому $\norm{\vect b}=d_{\min}$. Если же проецируется положение
в прежней TCA после манёвра, полученная величина приближённо
характеризует новый минимальный промах. Окончательный результат
уточняется повторным прогнозом и поиском нового минимума.

\coursefigure{concept-geometry}{Переход от локального относительного
движения к плоскости сближения. Слева показано сечение плоскостью,
содержащей относительную скорость и вектор промаха; справа --- две
координаты промаха. Радиус $R_c$ характеризует геометрическую
область столкновения.}

\subsection{Размеры объектов и длительность пролёта}

В простейшей геометрической модели каждый объект заменяется
сферой с радиусом, охватывающим его размеры. Сферы пересекаются,
если расстояние между центрами меньше суммы радиусов. Поэтому
один объект можно считать точкой, а второй --- сферой
с приведённым радиусом столкновения
\begin{equation}
R_c=R_{c,1}+R_{c,2}.
\label{eq:combined-radius}
\end{equation}
В плоскости сближения опасная область представляет собой круг
$\norm{\vect b}\leqslant R_c$. Для номинальных траекторий это
геометрический признак столкновения. Для случайного положения
необходимо определить вероятность попадания в данный круг.

Прямолинейное приближение должно быть верно на времени прохождения
не только размеров тел, но и существенной области неопределённости.
Если её характерный размер вдоль относительного движения равен
$L$, продолжительность пролёта имеет порядок
$L/\norm{\vect v_{\mathrm{rel}}}$. За это время относительная
траектория должна оставаться почти прямой, а ковариация положения
--- практически постоянной. При малой относительной скорости
время возрастает; кривизна траекторий и изменение ковариации
могут стать существенными~\cite{li,nasa}.

\section{Неопределённость и вероятность столкновения}
\label{sec:probability}

Прогнозируемое положение известно с ошибкой. Поэтому номинальные
траектории дают только одно ожидаемое сближение из множества
возможных реализаций. Вероятностная модель позволяет оценить,
какая доля этих реализаций приводит к пересечению тел.

\subsection{Ковариация относительного положения}

Обозначим ошибки инерциального положения в одной эпохе через
$\delta\vect R_1$ и $\delta\vect R_2$ и будем считать их средние
нулевыми. Ковариационная матрица первого объекта определяется как
\[
C_1=\mathbb{E}\bigl[\delta\vect R_1\delta\vect R_1^\top\bigr].
\]
Её диагональные элементы имеют размерность квадрата длины
и равны дисперсиям координат. Внедиагональные элементы описывают
связь ошибок разных координат. Собственные векторы определяют
главные направления эллипсоида неопределённости, а квадратные
корни собственных значений --- среднеквадратические ошибки
в этих направлениях.

Если ковариация задана в другом ортонормированном базисе,
она преобразуется по правилу $C_{\text{ИСК}}=QCQ^\top$,
где столбцы $Q$ содержат оси исходного базиса в инерциальных
координатах. Если ковариация относится к более ранней эпохе,
перед сравнением её требуется перенести к сближению.
Таким образом, сложение матриц имеет смысл только после
согласования эпохи, базиса и единиц измерения.

Для линейного переноса ковариации можно использовать ту же
матрицу изохронных производных, что и для приращений состояния.
Если $C_y(t_i)$ --- ковариация всех шести элементов,
а ошибками модели движения на интервале пренебрегаем, то из
$\delta\vect y(t)=\Phi(t,t_i)\delta\vect y(t_i)$ следует
\begin{equation}
 C_y(t)=\Phi(t,t_i)C_y(t_i)\Phi^\top(t,t_i),\qquad
 C_k(t)=\frac{\partial\vect R_k}{\partial\vect y_k}
 C_y(t)\left(\frac{\partial\vect R_k}{\partial\vect y_k}\right)^\top.
 \label{eq:covariance-propagation}
\end{equation}
Вторая формула переводит ковариацию элементов в ковариацию
положения объекта $k$ в общей инерциальной системе.
Перенос выполняется для каждого объекта отдельно.
Например, неопределённость большой полуоси приводит
к росту ошибки фазы при увеличении времени прогноза.

Ошибка относительного положения равна
$\delta\vect\rho=\delta\vect R_2-\delta\vect R_1$.
Для независимых ошибок её ковариация имеет вид
\begin{equation}
C_{\rho}=C_1+C_2.
\label{eq:combined-covariance}
\end{equation}
Если оценки взаимно коррелированы, требуется сохранить
перекрёстные слагаемые:
\[
C_{\rho}=C_1+C_2-C_{12}-C_{12}^\top,
\qquad
C_{12}=\mathbb{E}[\delta\vect R_1\delta\vect R_2^\top].
\]
Далее рассматривается независимость оценок объектов.

Линейная проекция случайного вектора преобразует ковариацию
тем же правилом, что и поворот координат. Поэтому
\begin{equation}
C_b=P C_{\rho}P^\top.
\label{eq:projected-covariance}
\end{equation}
Это матрица размера $2\times2$, описывающая разброс в плоскости
сближения. При гауссовом распределении контуры постоянной плотности
являются эллипсами
\begin{equation}
(\vect z-\vect b)^\top C_b^{-1}(\vect z-\vect b)=k^2.
\label{eq:covariance-ellipse}
\end{equation}
Их полуоси равны $k\sigma_1$ и $k\sigma_2$, где
$\sigma_1^2,\sigma_2^2$ --- собственные значения $C_b$.
Внутри такого двумерного эллипса заключена вероятность
$1-\exp(-k^2/2)$. Например, эллипс с полуосями
$\sigma_1,\sigma_2$ содержит около $39{,}3\%$ вероятности;
одномерное правило $68{,}3\%$ к нему не применяется.

\subsection{Вероятность за короткий пролёт}

Для быстрого сближения примем почти постоянную относительную
скорость и ковариацию за время пролёта. Влиянием неопределённости относительной скорости на геометрию
за время короткого пролёта пренебрегаем. Распределение ошибки
относительного положения считаем гауссовым.
В координатах $(z_1,z_2,s)$, где $s$ направлено вдоль относительной
скорости, опасные взаимные положения образуют цилиндр радиуса
$R_c$: продольная координата определяет, когда точка пройдёт
рядом с первым объектом, а $(z_1,z_2)$ --- насколько близко.

Пусть $p_3(\vect z,s)$ --- трёхмерная плотность относительного
положения в этих координатах. Вероятность столкновения
в прямолинейной модели записывается как
\[
P_c=\iint_{\norm{\vect z}\leqslant R_c}
       \left(\int_{-\infty}^{\infty}p_3(\vect z,s)\,\dd s\right)
       \dd z_1\dd z_2.
\]
Внутренний интеграл является двумерной маргинальной плотностью.
Такое интегрирование допустимо и при корреляции продольной
и поперечных координат; независимость этих координат не требуется.
Для гауссова распределения среднее проекции равно $\vect b$,
а ковариация --- $C_b$. В результате
\begin{equation}
\boxed{
P_c=\frac{1}{2\pi\sqrt{\det C_b}}
\iint_{\norm{\vect z}\leqslant R_c}
\exp\!\left[-\frac12(\vect z-\vect b)^\top
C_b^{-1}(\vect z-\vect b)\right]\dd z_1\dd z_2.}
\label{eq:pc-integral}
\end{equation}
Здесь круг столкновения расположен в начале координат,
а центр плотности --- в точке $\vect b$. Эквивалентное изображение
получается переносом начала координат в центр плотности:
тогда на вектор промаха смещается круг столкновения.

Формула~\eqref{eq:pc-integral} оценивает вероятность столкновения
\emph{за короткий пролёт}. Вероятность нахождения относительного
положения внутри сферы $\norm{\vect\rho}\leqslant R_c$
в единственный фиксированный момент --- другая величина.
При длительном сближении или существенном изменении распределения
ошибки двумерное сведение требует замены моделью,
учитывающей движение во времени~\cite{li,nasa}.

Для вычисления интеграла удобно повернуть оси плоскости вдоль
главных осей $C_b$. Круг столкновения при повороте не изменяется.
Пусть в новых координатах среднее равно $(\beta_1,\beta_2)^\top$,
а ковариация --- $\operatorname{diag}(\sigma_1^2,\sigma_2^2)$.
При фиксированной первой координате $x$ допустимая вторая
координата лежит между $-\sqrt{R_c^2-x^2}$ и
$\sqrt{R_c^2-x^2}$. Аналитическое интегрирование по ней даёт
\begin{equation}
\begin{split}
P_c={}&\frac{1}{2\sqrt{2\pi}\sigma_1}
\int_{-R_c}^{R_c}
\exp\!\left[-\frac{(x-\beta_1)^2}{2\sigma_1^2}\right]\\[-1mm]
&\quad\times\left[
\operatorname{erf}\!\left(
\frac{\sqrt{R_c^2-x^2}-\beta_2}{\sqrt2\sigma_2}\right)
-\operatorname{erf}\!\left(
\frac{-\sqrt{R_c^2-x^2}-\beta_2}{\sqrt2\sigma_2}\right)
\right]\dd x.
\end{split}
\label{eq:pc-1d}
\end{equation}
Здесь $\operatorname{erf}(q)=2\pi^{-1/2}\int_0^q\exp(-u^2)\dd u$.
Таким образом, для каждого варианта манёвра достаточно вычислить
одномерный интеграл. Для очень малых вероятностей разность
близких функций ошибок следует вычислять устойчивым способом,
например через дополнительные функции ошибок. Основные
схемы квадратур и условия их применения рассмотрены
в~\cite{nasa,li}.

\subsection{Геометрический промах и учёт ковариации}

Номинальное расстояние измеряет длину вектора $\vect b$,
но не учитывает направление неопределённости. На
рисунке~\ref{fig:concept-risk} в обоих случаях приняты
$\sigma_1=120$~м, $\sigma_2=30$~м и $R_c=10$~м.
При $\vect b=(0,80)^\top$~м получено
$P_c\approx4{,}30\cdot10^{-4}$, а при
$\vect b=(140,0)^\top$~м --- $P_c\approx6{,}94\cdot10^{-3}$.
Во втором случае номинальное расстояние больше, однако круг
столкновения находится в области большей плотности распределения.
Обе вероятности вычислены по~\eqref{eq:pc-1d}.

\coursefigure{concept-risk}{Два промаха при одной ковариационной
матрице. Эллипсы соответствуют $k=1$ и $k=2$ в
формуле~\eqref{eq:covariance-ellipse}; оранжевый круг --- область
столкновения. Увеличение евклидова расстояния с 80 до 140~м
в данном сравнении сопровождается ростом $P_c$.}

Часто для учёта неопределённости используется квадрат расстояния
Махаланобиса
\begin{equation}
D^2=\vect b^\top C_b^{-1}\vect b.
\label{eq:mahalanobis-distance}
\end{equation}
Он измеряет удаление в масштабе среднеквадратических ошибок.
Связь этого критерия с вероятностью и его применение
для выбора манёвра будут рассмотрены в разделе~\ref{sec:design}.

\subsection{Рост неопределённости и уменьшение вероятности}

Уменьшение $P_c$ может происходить по разным причинам. Если
манёвр удаляет центр распределения от области столкновения
при той же ковариации, оно характеризует результат управления.
Если же увеличивается сама неопределённость, вероятность
попадания в небольшой круг также может уменьшиться вследствие
распределения вероятности по большей области.

Для нулевого номинального промаха и изотропной ковариации
$C_b=\sigma^2I_2$ интеграл в полярных координатах вычисляется
непосредственно:
\begin{equation}
P_c=\int_0^{R_c}\frac{r}{\sigma^2}
\exp\!\left(-\frac{r^2}{2\sigma^2}\right)\dd r
=1-\exp\!\left(-\frac{R_c^2}{2\sigma^2}\right).
\label{eq:zero-miss-isotropic}
\end{equation}
При увеличении $\sigma$ это значение стремится к нулю,
хотя номинальные траектории по-прежнему пересекаются.

При ненулевом промахе зависимость обычно немонотонна.
На рисунке~\ref{fig:concept-dilution} расстояние равно 100~м,
радиус столкновения --- 10~м. При малой $\sigma$ положение
известно достаточно точно, и распределение почти не достигает
опасной области. При росте $\sigma$ вероятность увеличивается,
достигает максимума около $\sigma=70{,}5$~м,
а затем уменьшается. На последнем участке оценки положения
становятся менее определёнными.

\coursefigure[.88]{concept-dilution}{Вероятность столкновения
при фиксированном номинальном промахе 100~м, $R_c=10$~м
и ковариации $C_b=\sigma^2I_2$. Кривая рассчитана по
формуле~\eqref{eq:pc-1d}; максимум равен примерно
$3{,}68\cdot10^{-3}$.}

Этот эффект называется \emph{размыванием вероятности}
(probability dilution)~\cite{nasa}. Малое значение $P_c$
при большой неопределённости не означает, что объекты
с высокой точностью прогнозируются далеко друг от друга.
При принятии решения поэтому оценивают также достоверность
ковариаций и возможность уточнения орбит.

\begin{explanation}{Результат вероятностной оценки}
Опасность короткого сближения определяется совместно вектором
промаха, ковариацией и размерами объектов. Манёвр меняет геометрию
сближения; ошибки его исполнения могут изменить и ковариацию.
Поэтому при выборе управления необходимо переносить к событию
и номинальное воздействие импульса, и относящиеся к нему
неопределённости.
\end{explanation}


\section{Влияние импульсного манёвра на сближение}
\label{sec:influence}

Вероятность столкновения определяется положением среднего вектора
промаха относительно области столкновения и распределением ошибок.
Чтобы выбирать манёвр, требуется связать его параметры с изменением
этого вектора. Для малого импульса используем уже знакомую
матрицу манёвра в орбитальных элементах, а затем перейдём
от изменения элементов к смещению в плоскости сближения.

\subsection{От изменения скорости к изменению элементов}

Пусть $\Delta\vect V=(\Delta V_r,\Delta V_t,\Delta V_n)^\top$
--- компоненты импульса в орбитальной системе управляемого КА
непосредственно перед манёвром. Базис задаётся текущим состоянием:
\begin{equation}
 \vect e_r=\frac{\vect R_1}{R_1},\qquad
 \vect e_n=\frac{\vect R_1\times\vect V_1}
 {\norm{\vect R_1\times\vect V_1}},\qquad
 \vect e_t=\vect e_n\times\vect e_r,
 \qquad Q=[\vect e_r\ \vect e_t\ \vect e_n].
 \label{eq:rtn-basis}
\end{equation}
Трансверсальное направление ортогонально радиусу в плоскости орбиты;
на эллиптической орбите оно в общем случае отличается от направления
скорости. Точное условие идеального импульса имеет вид
\begin{equation}
 \vect R_1^+=\vect R_1^-,\qquad
 \vect V_1^+=\vect V_1^-+Q\Delta\vect V.
 \label{eq:impulse-jump}
\end{equation}
Элементы после манёвра вычисляются по этому положению и новой скорости.

Переход к линейной модели означает дифференцирование этого
преобразования при нулевом импульсе:
\begin{equation}
 \delta\vect y(t_m^+)=B(\vect y_1(t_m))\Delta\vect V,
 \qquad
 B=\left.
 \frac{\partial\vect y(\vect R_1,\vect V_1+Q\Delta\vect V)}
 {\partial\Delta\vect V}\right|_{\Delta\vect V=0}.
 \label{eq:B-definition}
\end{equation}
Здесь и далее линейные равенства описывают первый порядок
по малому управлению. Матрица $B$ совпадает с матрицей коэффициентов
при ускорениях в уравнениях Гаусса, используемой в задаче встречи.
Например, её первая строка даёт
\begin{equation}
 \delta a=\frac{2a^2}{\sqrt{\mu p}}
 \left(e\sin\nu\,\Delta V_r+(1+e\cos\nu)\Delta V_t\right),
 \qquad p=a(1-e^2).
 \label{eq:delta-a}
\end{equation}
Радиальный и трансверсальный импульсы меняют энергию и среднее
движение; нормальный импульс в первом порядке меняет ориентацию
плоскости орбиты.

\subsection{Перенос изменения к моменту сближения}

Обозначим через $\vect y_1(t)$ номинальное движение управляемого КА.
Матрица изохронных производных $\Phi(t,t_m)$ переносит малое
изменение элементов между заданными эпохами. В каждой эпохе
номинальное и возмущённое состояния сравниваются при одном
физическом времени.
Поэтому при номинальном моменте сближения
\begin{equation}
 \delta\vect y(t_{\mathrm{ca}})=M(t_m)\Delta\vect V,
 \qquad
 M(t_m)=\Phi(t_{\mathrm{ca}},t_m)B(\vect y_1(t_m)).
 \label{eq:M-definition}
\end{equation}
Все матрицы в \eqref{eq:M-definition} вычисляются вдоль одной
номинальной траектории управляемого аппарата. Движение опасного
объекта прогнозируется отдельно.

В центральном поле первые пять элементов постоянны,
а $\dot M=n(a)=\sqrt{\mu/a^3}$. В матрице $\Phi$ все
диагональные элементы равны единице; единственный ненулевой
внедиагональный элемент имеет индексы $(6,1)$:
\begin{equation}
 \Phi_{61}(t,t_m)=-\frac{3n}{2a}(t-t_m).
 \label{eq:kepler-phi}
\end{equation}
Отсюда следует физически важное соотношение
\begin{equation}
 \delta M(t_{\mathrm{ca}})=\delta M(t_m^+)
   -\frac{3n}{2a}(t_{\mathrm{ca}}-t_m)\delta a.
 \label{eq:phase-drift}
\end{equation}
Изменение большой полуоси создаёт накапливающийся фазовый сдвиг.
Именно поэтому малый заблаговременный импульс может дать заметное
смещение в момент сближения. Накопление фазы не означает,
что каждое направление импульса одинаково эффективно при любом
времени упреждения.

Положение определяется элементами, поэтому
\begin{equation}
 \delta\vect R_1(t_{\mathrm{ca}})=
 \frac{\partial\vect R_1}{\partial\vect y_1}
 M(t_m)\Delta\vect V.
 \label{eq:position-change}
\end{equation}
Производная здесь относится к текущим элементам в момент
$t_{\mathrm{ca}}$. Она не включает повторный перенос по времени:
этот перенос уже учтён в $M$.

Для вычисления производной можно продифференцировать обычное
восстановление положения по элементам. В собственной плоскости
орбиты
\begin{equation}
 \vect R_{\mathrm{orb}}=
 \begin{pmatrix}a(\cos E-e)\\
 a\sqrt{1-e^2}\sin E\\0\end{pmatrix},\qquad
 M=E-e\sin E,
 \label{eq:position-kepler}
\end{equation}
после чего выполняются повороты на $\omega,i,\Omega$.
При дифференцировании уравнения Кеплера полезны формулы
\begin{equation}
 \left.\frac{\partial E}{\partial M}\right|_e
 =\frac{1}{1-e\cos E},\qquad
 \left.\frac{\partial E}{\partial e}\right|_M
 =\frac{\sin E}{1-e\cos E}.
 \label{eq:kepler-derivatives}
\end{equation}
Альтернативный вычислительный способ --- центральные разности
положения по каждому элементу с сохранением общей эпохи.

\subsection{Матрица смещения в плоскости сближения}

В принятом определении $\vect\rho=\vect R_2-\vect R_1$
манёвр первого аппарата даёт $\delta\vect\rho=-\delta\vect R_1$.
Зафиксируем плоскость номинального сближения и её матрицу $P$.
Тогда из \eqref{eq:position-change} получаем
\begin{equation}
 \boxed{\quad
 \vect b(\Delta\vect V)\simeq\vect b_0+G(t_m)\Delta\vect V,
 \qquad
 G(t_m)=-P\frac{\partial\vect R_1}{\partial\vect y_1}M(t_m).
 \quad}
 \label{eq:G-definition}
\end{equation}
Матрица $G$ имеет размер $2\times3$.
Её столбцы показывают смещение вектора промаха от единичных
радиального, трансверсального и нормального импульсов.
Если длина измеряется в метрах, а импульс в метрах в секунду,
элементы $G$ имеют размерность секунды.

Ту же матрицу можно вычислить непосредственно по прогнозам
с малыми пробными импульсами. Пусть $\vect d_j$ --- единичное
направление одной из осей ОСК, а $\varepsilon$ --- величина
пробного импульса. При фиксированных $P$ и $t_{\mathrm{ca}}$
\begin{equation}
 G_{:,j}\simeq
 \frac{\vect b(+\varepsilon\vect d_j)
       -\vect b(-\varepsilon\vect d_j)}{2\varepsilon}.
 \label{eq:G-finite-difference}
\end{equation}
Каждый прогноз использует точный скачок \eqref{eq:impulse-jump}.
Так можно получить чувствительность и для возмущённой модели,
не выводя аналитическую матрицу перехода. Уменьшение шага
проверяется сравнением результатов нескольких разностей;
слишком малый шаг усиливает влияние ошибок вычисления положения.

\begin{explanation}{Что именно линеаризуется}
Малым считается отклонение управляемого КА от его собственной
номинальной траектории. Опасный объект может двигаться по существенно
другой орбите и приближаться с большой относительной скоростью.
Уравнения CW для двух близких орбит не требуется применять
к его движению на всём интервале до сближения.
\end{explanation}

\subsection{Изменение времени наибольшего сближения}

Манёвр меняет не только положение, но и время события. Для
невырожденного локального минимума расстояния имеем
$\vect\rho\cdot\vect v_{\mathrm{rel}}=0$ и
$v_{\mathrm{rel}}^2+\vect\rho\cdot\vect a_{\mathrm{rel}}>0$.
Линеаризация первого условия с учётом изменения времени даёт
\begin{equation}
 \delta t_{\mathrm{ca}}=-
 \frac{\vect v_{\mathrm{rel}}\cdot\delta\vect\rho
       +\vect\rho\cdot\delta\vect v_{\mathrm{rel}}}
 {v_{\mathrm{rel}}^2+\vect\rho\cdot\vect a_{\mathrm{rel}}},
 \qquad
 \vect a_{\mathrm{rel}}=\dot{\vect V}_2-\dot{\vect V}_1.
 \label{eq:event-time}
\end{equation}
Все величины в правой части относятся к номинальному
$t_{\mathrm{ca}}$. Дополнительное изменение положения из-за
сдвига эпохи равно $\vect v_{\mathrm{rel}}\delta t_{\mathrm{ca}}$.
Его проекция на номинальную плоскость исчезает:
\[
 P\vect v_{\mathrm{rel}}\delta t_{\mathrm{ca}}=0.
\]
Это объясняет, почему \eqref{eq:G-definition} описывает первый
порядок поперечного смещения без отдельного вычисления
$\delta t_{\mathrm{ca}}$.

Однако направление относительной скорости тоже меняется.
Новая плоскость и её ковариация отличаются от номинальных;
координаты в заново выбранном базисе нельзя буквально
отождествлять с прежними. Поэтому линейный синтез выполняется
с фиксированной плоскостью, а итоговая проверка заново определяет
$t_{\mathrm{ca}}$, $P$, $C_b$ и $P_c$.

\section{Выбор величины, направления и времени манёвра}
\label{sec:design}

Матрица $G$ позволяет быстро оценивать большое число малых
импульсов. Теперь необходимо выбрать критерий и учесть допустимые
параметры манёвра. Рассмотрим последовательно поиск при заданном
направлении, аналитический выбор направления для простейшей
геометрической задачи и минимизацию характеристической скорости
с ограничением вероятности.

\subsection{Поиск при заданном направлении}

Пусть момент $t_m$ и единичное направление $\vect d$ заданы,
а импульс равен $\Delta\vect V=s\vect d$, $s\geqslant0$.
В плоскости сближения его действие описывается прямой
\begin{equation}
 \vect b(s)=\vect b_0+sG(t_m)\vect d.
 \label{eq:fixed-direction}
\end{equation}
Изменяя $s$, вычисляем $P_c(\vect b(s),C_b,R_c)$ по интегралу
вероятности и находим первое допустимое значение.
Если разрешены оба знака импульса, направления $\vect d$
и $-\vect d$ рассматриваются отдельно.

Начальное движение вдоль прямой может приближать среднее
к области столкновения, поэтому вероятность не обязана
уменьшаться от $s=0$. Сначала следует просмотреть допустимый
интервал величины импульса, выделить переходы через
$P_{\text{доп}}$, а затем уточнить границы одномерным поиском
корня. Если направлением нельзя достаточно изменить вектор
промаха при заданном пределе $s$, допустимого решения в этом
семействе нет.

Повторение такого поиска по нескольким моментам и направлениям
даёт карту допустимых манёвров. Выбирается минимальный импульс
среди проверенных вариантов. Если проверены только оси ОСК,
это минимум среди осевых манёвров; произвольное пространственное
направление может потребовать меньшей характеристической скорости.

\subsection{Наибольший промах при заданном импульсе}

Сначала решим геометрическую задачу без ковариации: при заданном
моменте и величине $d$ выбрать направление наибольшего промаха.
При нулевом номинальном промахе $\vect b_0=0$
\begin{equation}
 \begin{aligned}
 \max_{\Delta\vect V}\quad
  &\Delta\vect V^\top G^\top G\Delta\vect V,\\
 \text{s.t.}\quad
  &\norm{\Delta\vect V}=d.
 \end{aligned}
 \label{eq:maximum-deviation}
\end{equation}
Критерий равен квадрату промаха $\norm{\vect b}^2$.
Матрица $G^\top G$ симметрична и неотрицательно определена.
В её ортонормированном собственном базисе квадратичная форма
равна сумме $\lambda_j v_j^2$. При условии
$\sum_jv_j^2=d^2$ наибольшее значение получается при размещении
всей нормы в собственном направлении с максимальным
собственным значением:
\begin{equation}
 \Delta\vect V=\pm d\vect d_{\max},\qquad
 G^\top G\vect d_{\max}=\lambda_{\max}\vect d_{\max},
 \qquad
 b_{\max}=d\sqrt{\lambda_{\max}}.
 \label{eq:principal-direction}
\end{equation}
Геометрически шар допустимых импульсов
отображается матрицей $G$ в заполненный эллипс смещений. Формула
\eqref{eq:principal-direction} выбирает его большую полуось.
Это частный аналитический результат, используемый в синтезе
манёвров уклонения~\cite{gonzalo}.

При ненулевом $\vect b_0$ постановка принимает вид
\begin{equation}
 \begin{aligned}
 \max_{\Delta\vect V}\quad
  &\norm{\vect b_0+G\Delta\vect V}^2,\\
 \text{s.t.}\quad
  &\norm{\Delta\vect V}=d.
 \end{aligned}
 \label{eq:affine-quadratic}
\end{equation}
Раскрывая квадрат нормы в критерии, получаем
\[
 \norm{\vect b_0+G\Delta\vect V}^2
 =\norm{\vect b_0}^2
 +2\Delta\vect V^\top G^\top\vect b_0
 +\Delta\vect V^\top G^\top G\Delta\vect V.
\]
Появляется линейный член; направления $\Delta\vect V$
и $-\Delta\vect V$ больше не равноценны.
Простой выбор главного собственного вектора в общем случае
не решает \eqref{eq:affine-quadratic}. В расчётном примере эта
двумерная по углам задача решается численно на сфере
$\norm{\Delta\vect V}=d$.

\subsection{Направление смещения с учётом неопределённости}

Большой евклидов промах не обязательно соответствует меньшей
вероятности столкновения. При анизотропной ковариации важны
положение среднего относительно короткой и длинной осей
эллипса и размер области столкновения. Полезный приближённый
критерий --- квадрат расстояния Махаланобиса:
\begin{equation}
 D^2=\vect b^\top C_b^{-1}\vect b.
 \label{eq:mahalanobis-design}
\end{equation}
После перехода к главным осям ковариации эта величина равна
$b_1^2/\sigma_1^2+b_2^2/\sigma_2^2$. Смещение на одинаковое
расстояние вдоль короткой оси сильнее меняет $D^2$.

Если плотность мало меняется внутри круга столкновения,
интеграл приближённо равен площади круга, умноженной на
плотность в его центре:
\begin{equation}
 P_c\simeq\frac{R_c^2}{2\sqrt{\det C_b}}
             \exp\left(-\frac{D^2}{2}\right).
 \label{eq:small-body-pc}
\end{equation}
При фиксированных $C_b$ и $R_c$ максимизация $D^2$ тогда
приближённо минимизирует $P_c$. Для $\vect b_0=0$
повторяется собственная задача \eqref{eq:principal-direction},
но с матрицей $G^\top C_b^{-1}G$ вместо $G^\top G$.
В общем случае соответствующие направления различаются.

Условие малости области относится к изменению плотности,
а не только к отношению радиуса к полуосям эллипса.
Действительно, для $\norm{\vect z}\leqslant R_c$
\begin{equation}
 \log\frac{f(\vect z)}{f(0)}
 =\vect z^\top C_b^{-1}\vect b
    -\frac12\vect z^\top C_b^{-1}\vect z.
 \label{eq:density-variation}
\end{equation}
Достаточно, чтобы обе величины
$R_c\norm{C_b^{-1}\vect b}$ и
$R_c^2\lambda_{\max}(C_b^{-1})$ были малы.
При удалённом среднем первая из них может быть существенной
даже при $R_c\ll\sigma_{\min}$.

В анизотропном случае преобразование, превращающее ковариацию
в единичную, одновременно превращает круг столкновения в эллипс.
Поэтому точный интеграл по кругу в общем случае не определяется
одним $D^2$. Приближения вероятности удобны для понимания
геометрии и получения начального направления; окончательное
ограничение в нашем расчёте проверяется непосредственной квадратурой.

\subsection{Минимальная характеристическая скорость}

Для фиксированного $t_m$ и фиксированной номинальной
ковариации основная задача в линейной модели имеет вид
\begin{equation}
 \begin{aligned}
 \min_{\mathclap{\Delta\vect V\in\mathcal A(t_m)}}\quad
  &\norm{\Delta\vect V},\\
 \text{s.t.}\quad
  &P_c(\vect b_0+G(t_m)\Delta\vect V,C_b,R_c)
      \leqslant P_{\text{доп}}.
 \end{aligned}
 \label{eq:minimum-dv}
\end{equation}
Множество $\mathcal A$ задаёт допустимые направления и величину
импульса. При свободном времени расчёт повторяется для
$t_m$ из разрешённого интервала либо время включается
в число переменных оптимизации.

Ограничение безопасности обычно требует находиться вне
области повышенной вероятности. При изотропной ковариации
внешность круга не выпукла: отрезок между двумя
безопасными точками может проходить через опасную область.
Следовательно, свойства выпуклой линейной задачи встречи
с равенством конечных состояний не переносятся автоматически
на \eqref{eq:minimum-dv}.

Практический расчёт можно построить так: получить допустимые
начальные приближения из одномерных поисков, уточнить все
три компоненты импульса с ограничением вероятности и сравнить
полученные решения. Несколько начальных направлений уменьшают
зависимость от одного локального решения. Они не являются
доказательством глобального минимума.

При необходимости используют последовательную линеаризацию
ограничения или последовательную выпуклую оптимизацию.
Каждая подзадача относится к локальному приближению исходной
задачи; окончательная допустимость проверяется исходной моделью.
Обзор импульсных и протяжённых методов, включая такие подходы,
приведён в~\cite{review}.

\begin{explanation}{Три разных результата расчёта}
Максимальный промах при заданной величине импульса,
минимальная вероятность при той же величине и минимальная
характеристическая скорость при заданной допустимой вероятности
--- разные задачи. При сравнении решений необходимо фиксировать,
какие величины заданы и какой критерий оптимизируется.
\end{explanation}

\section{Расчёт уклонения на пространственных орбитах}
\label{sec:example}

Рассмотрим один случай, в котором можно проследить всю
последовательность: построение геометрии, вычисление вероятности,
выбор направления и величины импульса, затем повторный прогноз.
Пример синтетический: его параметры выбраны для сопоставления
геометрического и вероятностного критериев. Это не восстановление
реального сближения по каталожным данным.

\subsection{Орбиты и номинальное сближение}

Используем центральное поле Земли с
$\mu=398600{,}4418$~км$^3$/с$^2$. Номинальную TCA примем
за начало отсчёта времени. Элементы управляемого аппарата в этой
эпохе:
\begin{equation}
 a_1=7000~\text{км},\quad e_1=0{,}01,\quad
 \omega_1=40^\circ,\quad i_1=51^\circ,\quad
 \Omega_1=20^\circ,\quad\nu_1=60^\circ.
 \label{eq:example-primary}
\end{equation}
Средняя аномалия, соответствующая заданной истинной аномалии,
равна примерно $59{,}0113^\circ$; в расчёте она вычисляется
через уравнение Кеплера без округления промежуточных значений.
Период составляет $T=\ExamplePeriod$~с.

Чтобы однозначно задать второй объект и малый промах,
построим его состояние непосредственно в эпоху $t=0$.
Радиальную компоненту скорости сохраним, а трансверсальную
повернём на $60^\circ$ вокруг радиального направления первого КА:
\begin{equation}
 \vect V_2=V_{r1}\vect e_{r1}
       +V_{t1}\cos60^\circ\,\vect e_{t1}
       +V_{t1}\sin60^\circ\,\vect e_{n1}.
 \label{eq:example-secondary-velocity}
\end{equation}
Относительная скорость перпендикулярна $\vect e_{r1}$.
Выберем $\vect e_1=\vect e_{r1}$,
$\vect e_2=\vect e_v\times\vect e_1$ и зададим
\begin{equation}
 \vect b_0=\begin{pmatrix}25\\-15\end{pmatrix}\text{м},
 \qquad
 \vect R_2=\vect R_1+P^\top\vect b_0.
 \label{eq:example-secondary-position}
\end{equation}
В последней формуле положения и промах предварительно
выражаются в одних единицах. По $(\vect R_2,\vect V_2)$
восстанавливается эллиптическая орбита второго объекта:
\[
a_2\approx7000{,}050512~\text{км},\qquad e_2\approx0{,}0100002244,
\]
\[
i_2\approx64{,}43534^\circ,\quad
\Omega_2\approx90{,}98391^\circ,\quad
\omega_2\approx-1{,}93879^\circ.
\]

По построению $\vect\rho\cdot\vect v_{\mathrm{rel}}=0$;
проверка второй производной подтверждает локальный минимум.
Номинальный промах равен $\ExampleNominalMiss$~м,
относительная скорость --- $\ExampleRelativeSpeed$~км/с.
Таким образом, близость объектов в одну эпоху сочетается
с существенно различными пространственными орбитами.

\coursefigure{example-orbits}{Номинальные орбиты двух объектов
и положение сближения. Участки до и после события принадлежат
разным орбитальным плоскостям. Расстояние между объектами
в TCA мало по сравнению с радиусом орбиты и в этом масштабе
неразличимо.}

\subsection{Ковариация и исходная вероятность}

Примем $R_c=10$~м и $P_{\text{доп}}=10^{-5}$.
Ковариация относительного положения задана непосредственно
в эпоху сближения. Её главные среднеквадратические ошибки
в плоскости составляют 40 и 200~м; ось, соответствующая
40~м, повёрнута на $35^\circ$ от $b_1$ к $b_2$.
Следовательно,
\begin{equation}
 C_b=\begin{pmatrix}
 14233{,}213&-18042{,}098\\
 -18042{,}098&27366{,}787
 \end{pmatrix}\text{м}^2.
 \label{eq:example-covariance}
\end{equation}
Для повторного проектирования после манёвра определим
трёхмерную ковариацию
\begin{equation}
 C_{\rho}=P^\top C_bP+(300~\text{м})^2\vect e_v\vect e_v^\top.
 \label{eq:example-covariance-3d}
\end{equation}
В короткой окрестности события эту матрицу будем считать
постоянной в инерциальных осях. Такой способ задаёт локальную
учебную модель; он не имитирует определение орбиты по измерениям.

Интегрирование \eqref{eq:pc-integral} даёт
\begin{equation}
 P_{c,0}=\ExampleNominalPc.
 \label{eq:example-initial-pc}
\end{equation}
Значение превышает выбранный допустимый уровень.
Характерный продольный размер $3\sigma=900$~м
проходится примерно за $0{,}12$~с. Орбитальная геометрия
за это время меняется мало; краткость пролёта согласуется
с применением двумерной модели.

\subsection{Сравнение направлений и критериев}

Основной момент манёвра выберем за $0{,}4T$ до события:
$t_m=-\ExampleLead$~с. Этот момент позволяет сравнить
три направления ОСК при ненулевом влиянии каждого из них.
Центральными пробными импульсами величины
$0{,}001$~м/с получена матрица
\begin{equation}
 G=\begin{pmatrix}
 -483{,}069&-3321{,}815&0\\
 2882{,}996&4276{,}692&-256{,}984
 \end{pmatrix}\text{с}.
 \label{eq:example-G}
\end{equation}
Практически нулевой элемент первой строки в нормальном
столбце следует из выбранной геометрии.

Для каждого осевого направления найдём наименьший по модулю
импульс обоих знаков, при котором линейная модель даёт
$P_c=10^{-5}$. Результаты приведены в
таблице~\ref{tab:axis}. Различие промахов на одной границе
вероятности отражает анизотропию ковариации.

\begin{table}[htbp]\centering
\caption{Осевые манёвры при $t_m=-0{,}4T$.
Величины и промахи определены в линейной модели
при одном ограничении $P_c=10^{-5}$.}
\label{tab:axis}
\begin{tabular}{@{}lrr@{}}\toprule
Направление&Компонента импульса, м/с&Промах, м\\\midrule
Радиальное&$0{,}099133$&$271{,}77$\\
Трансверсальное&$-0{,}121313$&$684{,}20$\\
Нормальное&$-0{,}876068$&$211{,}62$\\
\bottomrule\end{tabular}

\end{table}

Теперь разрешим произвольное направление и зафиксируем
$\norm{\Delta\vect V}=0{,}1$~м/с. Отдельно максимизируем
евклидов промах и минимизируем точный интеграл вероятности.
Поскольку исходный промах ненулевой, первая задача решается
с полным выражением \eqref{eq:affine-quadratic}.

\begin{table}[htbp]\centering
\caption{Сопоставление критериев при $t_m=-0{,}4T$.
Промах и вероятность относятся к линейной модели. В первых
двух строках модуль импульса одинаков; в последней он минимизируется
при $P_c\leqslant10^{-5}$.}
\label{tab:criteria}
\begin{tabular}{@{}lrrr@{}}\toprule
Задача&$\Delta V$, м/с&Промах, м&$P_c$\\\midrule
Максимальный промах&$0{,}100000$&$629{,}11$&$3{,}90\cdot10^{-5}$\\
Минимальная вероятность&$0{,}100000$&$358{,}00$&$7{,}63\cdot10^{-6}$\\
Минимальный импульс&$0{,}097840$&$348{,}63$&$1{,}00\cdot10^{-5}$\\
\bottomrule\end{tabular}

\end{table}

Максимизация промаха даёт около 629~м, но вероятность
$3{,}90\cdot10^{-5}$ остаётся выше допустимой. При той же
характеристической скорости другое направление даёт промах
около 358~м и вероятность $7{,}63\cdot10^{-6}$.
Во втором случае среднее удаляется в более благоприятном
направлении относительно короткой оси эллипса.
Вектор промаха и ковариация показаны на
рисунке~\ref{fig:example-bplane}.

\coursefigure{example-bplane}{Сравнение манёвров в плоскости
сближения. Цветные линии --- уровни $P_c$ как функции среднего
вектора $\vect b$; красный круг --- область столкновения
радиуса $R_c$. Пунктир задаёт масштаб неопределённости:
$\vect b^\top C_b^{-1}\vect b=1$. Максимальный промах
и минимальная вероятность при одинаковой норме импульса
достигаются разными направлениями.}

Минимизация нормы импульса с ограничением $P_c\leqslant10^{-5}$
из нескольких начальных приближений дала лучший найденный
вариант
\begin{equation}
 \Delta\vect V\approx
 \begin{pmatrix}0{,}095635\\0{,}017739\\-0{,}010576\end{pmatrix}
 \text{м/с},\qquad
 J=\ExampleOptimalDV~\text{м/с}.
 \label{eq:example-optimal-impulse}
\end{equation}
Это результат для выбранного времени и ограничений модели.
Наличие нескольких стартов не доказывает глобальную
оптимальность по всем возможным управлениям.

\subsection{Время упреждения и ограничение направления}

При изменении $t_m$ заново вычисляются ОСК в момент манёвра
и матрица $G$. На рисунке~\ref{fig:example-lead-time}
показаны чувствительности и требуемые величины осевых манёвров.
Точки свободного выбора направления соответствуют
отдельным решениям задачи \eqref{eq:minimum-dv}.
Ковариация в этом сравнении одна и та же и задана в TCA;
изменяется только время приложения управления.

\coursefigure{example-lead-time}{Влияние упреждения при
фиксированной ковариации в момент сближения. Чувствительности
показывают смещение от единичных импульсов; требуемая
характеристическая скорость относится к условию
$P_c\leqslant10^{-5}$. Отсутствие решения в показанном
диапазоне осевых импульсов не означает отсутствия
пространственного манёвра.}

За четверть периода до события лучший найденный свободный
импульс составляет около $0{,}10605$~м/с, за половину периода
--- $0{,}04895$~м/с, за один период --- $0{,}01453$~м/с,
за два периода --- $0{,}007265$~м/с.
На целых периодах особенно эффективно накапливается
фазовый сдвиг от изменения большой полуоси. При этом некоторые
другие столбцы $G$ могут быть малы, и соответствующий
осевой манёвр оказывается неэффективным.

Для одного трансверсального направления полную картину
даёт карта на рисунке~\ref{fig:example-risk-map}.
Она показывает знак и величину импульса, время его приложения
и границу допустимой вероятности. Такая карта позволяет
выбрать манёвр с учётом окна подготовки и доступного
направления тяги, прежде чем уточнять его параметры.

\coursefigure{example-risk-map}{Вероятность после
трансверсального импульса при разных величинах и временах
упреждения. Контур $P_c=10^{-5}$ отделяет допустимые варианты
в принятой линейной модели; цветовая шкала логарифмическая.}

Выигрыш раннего манёвра здесь показан при фиксированной
неопределённости в TCA. Если исходные измерения относятся
к разным эпохам или меняется срок прогноза ковариации,
сравнение следует выполнять заново. В операциях увеличение
упреждения также сопоставляется с возможностью получить
уточнённые данные до принятия решения.

\subsection{Повторный прогноз и численная проверка}

Для \eqref{eq:example-optimal-impulse} заново восстановим
элементы после точного скачка скорости и распространим обе
орбиты. Пассивное движение на всех участках вычисляется
в классических элементах: изменяется средняя аномалия,
затем решается уравнение Кеплера. Линейные приращения
элементов при этом не используются.

Поиск корня $\vect\rho\cdot\vect v_{\mathrm{rel}}=0$
в окрестности события даёт сдвиг TCA
$\ExampleOptimalTCA$~с. Новый промах равен
$\ExampleOptimalNonlinearMiss$~м вместо линейного
$\ExampleOptimalMiss$~м. После пересоздания плоскости
и проектирования \eqref{eq:example-covariance-3d}
получается $P_c\approx9{,}9992\cdot10^{-6}$.
Близость результатов подтверждает точность линейного
предсказания для данного малого импульса. Запас
относительно порога здесь практически отсутствует.

\coursefigure{example-local-encounter}{Расстояние в короткой
окрестности события до и после манёвра. Моменты минимума
определяются отдельно. Рост расстояния на долях секунды
обусловлен большой относительной скоростью. Справа ---
исходный и управляемый векторы промаха в номинальной плоскости;
пунктир --- эллипс $\vect b^\top C_b^{-1}\vect b=1$,
красный круг --- область столкновения. Показан номинальный
импульс~\eqref{eq:example-optimal-impulse}.}

Интеграл вероятности вычисляется квадратурой Гаусса
после замены $x=R_c\sin\theta$ в \eqref{eq:pc-1d}.
Сравнение 64 и 128 узлов и независимая адаптивная квадратура
проверяют численное интегрирование. Для матрицы $G$
сопоставлены шаги $0{,}0005$, $0{,}001$ и $0{,}002$~м/с;
относительное изменение составляет порядка $10^{-9}$.
Эти проверки относятся к вычислительной точности принятой
модели. Они не оценивают достоверность реальных измерений
и не заменяют добавление необходимых возмущающих сил.

\section{Ошибки исполнения и проверка манёвра}
\label{sec:verification}

Синтезированный импульс рассчитан по ожидаемому воздействию
двигательной установки. Реальное исполнение имеет ошибки,
а изменение орбиты может создать другие нежелательные события.
Поэтому итоговый расчёт должен учитывать не только номинальную
вероятность исходного сближения.

\subsection{Перенос ошибки импульса к сближению}

Пусть фактически выполненный импульс равен
$\Delta\vect V+\vect\varepsilon_V$, где среднее ошибки
нулевое, а ковариация равна $C_{\Delta V}$. В линейной модели
добавочная ошибка промаха составляет $G\vect\varepsilon_V$.
Если она независима от исходных ошибок состояния,
\begin{equation}
 C_{b,\mathrm{tot}}=C_b+G C_{\Delta V}G^\top.
 \label{eq:execution-covariance}
\end{equation}
Здесь добавляется только та неопределённость, которая ещё
не учтена в исходной ковариации. Ненулевое среднее ошибки
сначала смещает ожидаемый вектор промаха; одна добавка
к дисперсии не описывает систематическую ошибку.

Для примера примем независимые ошибки трёх компонент
с одинаковым среднеквадратическим значением
$\sigma_V=0{,}005$~м/с:
\begin{equation}
 C_{\Delta V}=\sigma_V^2 I_3.
 \label{eq:example-execution-error}
\end{equation}
Эта абсолютная модель выбрана как отдельное условие задачи.
Она не является характеристикой конкретной двигательной установки.
Для импульса \eqref{eq:example-optimal-impulse} вероятность
с учётом \eqref{eq:execution-covariance} возрастает до
$\ExampleExecutionPc$, то есть превышает $10^{-5}$.
Следовательно, номинальный расчёт на границе допустимости
не оставляет запаса для принятой ошибки исполнения.

Повторим оптимизацию с $C_{b,\mathrm{tot}}$ и потребуем
расчётное значение не более $8\cdot10^{-6}$.
Это выбранный запас относительно исходного порога.
Полученный импульс и результат его повторного прогноза
приведены ниже:
\begin{equation}
\Delta\vect V=\begin{pmatrix}0{,}098719\\0{,}018319\\-0{,}010917\end{pmatrix}\ \text{м/с},\qquad \norm{\Delta\vect V}=\ExampleSafeDV\ \text{м/с}.
\label{eq:example-safe-impulse}\end{equation}
\begin{equation}\begin{aligned}
\Delta t_{\mathrm{TCA}}&=\ExampleSafeTCA\ \text{с},\\
d_{\min}&=\ExampleSafeNonlinearMiss\ \text{м},\\
P_c\ \text{с учётом ошибок исполнения}&=\ExampleSafeNonlinearPc<10^{-5}.
\end{aligned}\label{eq:example-safe-verification}\end{equation}

В повторе ковариация исходного положения проектируется
на новую плоскость, а чувствительность к ошибке импульса
пересчитывается около найденного управления. Поэтому
проверяется не только старая матрица $G$, но и локальный
отклик уточнённой траектории.

\subsection{Протяжённый манёвр и ограничения миссии}

В модели конечной тяги приращение элементов определяется
воздействием на всём времени манёвра. Влияние протяжённого
манёвра на промах в первом порядке определяется интегралом
\begin{equation}
 \delta\vect b=-P\frac{\partial\vect R_1}{\partial\vect y_1}
 \int_{t_b}^{t_e}\Phi(t_{\mathrm{ca}},t)B(\vect y_1(t))
       \vect a(t)\,\dd t.
 \label{eq:finite-thrust-b}
\end{equation}
Величины $t_b,t_e$ означают начало и конец манёвра,
как в предыдущей лекции. При постоянной тяге ускорение
меняется вследствие расхода массы; закон направления
также должен быть задан.

Возможность заменить \eqref{eq:finite-thrust-b} импульсом
определяется изменением матрицы влияния за время манёвра.
Если оно мало, центрирование относительно $t_m$ уменьшает
ошибку импульсного приближения. При продолжительном манёвре
или малом оставшемся времени до сближения требуется
непосредственный прогноз с конечной тягой.

После уклонения меняются элементы рабочей орбиты.
Для найденного импульса~\eqref{eq:example-safe-impulse}
в центральном поле получены изменения
\[
 \delta a\approx32{,}18~\text{м},\qquad
 \delta e\approx-1{,}270\cdot10^{-5},\qquad
 \delta i\approx-5{,}74\cdot10^{-5}\text{ град}.
\]
Если эти изменения выходят за допуски миссии, восстановление
орбиты необходимо включить в дальнейшую последовательность
управления и в оценку суммарной характеристической скорости.
Восстановительный манёвр сам также проверяется на сближения.

\subsection{Последовательность итоговой проверки}

Сначала рассчитывают полную траекторию с найденным управлением
и принятой моделью возмущений. В окрестности исходного события
заново ищут TCA, определяют относительную скорость,
перестраивают плоскость и вычисляют вероятность с согласованной
ковариацией. Затем проверяют остальные ограничения:
время подготовки, доступное направление, ресурс ДУ,
параметры рабочей орбиты и другие опасные сближения
на установленном интервале прогноза.

Если второй объект тоже управляемый, прогноз должен учитывать
согласованный план его действий. Независимые изменения обеих
траекторий могут обесценить рассчитанный вариант.
При поступлении новых измерений исходные состояния и ковариации
обновляются; актуальность манёвра оценивается повторно.
После исполнения прогноз продолжается от достигнутого состояния.
Такой цикл связывает расчёт манёвра с практикой оценки
и снижения опасности сближений~\cite{cara}.

В нашем примере проверены исходное парное сближение,
нелинейность движения и заданная модель ошибки импульса.
Проверка по каталогу других объектов и ограничения конкретной
миссии в учебные исходные данные не входят. Эти части требуют
дополнительных данных, а не следуют из малого значения
вероятности одного события.

\subsection{Воспроизведение расчёта}

Расчётные данные сохранены с полной рабочей точностью;
приведённые в тексте числа округлены для чтения.
Скрипт \texttt{images/avoidance\_example.py} содержит
исходные состояния, преобразования элементов, квадратуру,
параметры оптимизации и проверочные прогнозы.
Он сохраняет результаты в \texttt{example-results.json}
и массивы для графиков в \texttt{example-data.npz}.
Концептуальные рисунки и их данные воспроизводит
\texttt{images/geometry\_probability.py}.
Из корня курса выполняются команды:
\begin{quote}\small\ttfamily
python3 9/images/geometry\_probability.py\\
python3 9/images/avoidance\_example.py\\
latexmk 9/collision-avoidance.tex
\end{quote}

\FloatBarrier
\begin{thebibliography}{9}
\addcontentsline{toc}{section}{Литература}
\bibitem{gonzalo}
Gonzalo J.~L., Colombo C., Di~Lizia P.
Analytical framework for space debris collision avoidance maneuver design.
\emph{Journal of Guidance, Control, and Dynamics}. 2021.
Vol.~44, no.~3. P.~469--487.
\href{https://doi.org/10.2514/1.G005398}{DOI: 10.2514/1.G005398}.
Препринт: \href{https://arxiv.org/abs/2012.12651}{arXiv:2012.12651}.
Геометрия и аналитический синтез: разделы II.A--D.

\bibitem{nasa}
NASA. \emph{NASA Spacecraft Conjunction Assessment and Collision
Avoidance Best Practices Handbook}. NASA/SP-20230002470, Rev.~1.
2023. Приложения M--N: вероятность столкновения,
размывание вероятности и вычислительные модели.
\href{https://ntrs.nasa.gov/citations/20230002470}{NASA NTRS}.

\bibitem{li}
Li J.-S., Yang Z., Luo Y.-Z.
A review of space-object collision probability computation methods.
\emph{Astrodynamics}. 2022. Vol.~6, no.~2. P.~95--120.
\href{https://doi.org/10.1007/s42064-021-0125-x}{DOI: 10.1007/s42064-021-0125-x}.

\bibitem{review}
Pavanello Z., De~Maria L., De~Vittori A., Maestrini M.,
Di~Lizia P., Armellin R.
CAMmary: A Review of Spacecraft Collision Avoidance Manoeuvre Design Methods.
2025. \href{https://arxiv.org/abs/2503.22555}{arXiv:2503.22555}.
Численные методы: раздел 3; аналитические методы: раздел 4.

\bibitem{cara}
NASA CARA. \emph{Step 3: Close Approach Risk Mitigation}.
\url{https://www.nasa.gov/cara/step-3-close-approach-risk-mitigation/}.
Дата обращения: 13 ноября 2026 г.
\end{thebibliography}
\end{document}
